% ============================================================
% Close-up: GoogleSearcher — Workflow Interface View
% Fragment — % ============================================================
% Close-up: GoogleSearcher — Workflow Interface View
% Fragment — % ============================================================
% Close-up: GoogleSearcher — Workflow Interface View
% Fragment — % ============================================================
% Close-up: GoogleSearcher — Workflow Interface View
% Fragment — \input{figs/closeup_google}
% ============================================================
\begin{figure}[htbp]
\centering
\resizebox{\textwidth}{!}{%
\begin{tikzpicture}[node distance=0.35cm, font=\sffamily,
    stepbar/.style={rectangle, rounded corners=6pt, thick,
        minimum width=19cm, align=center, font=\sffamily,
        drop shadow={opacity=0.10}, inner sep=6pt, text width=18cm},
    rwpanel/.style={rectangle, rounded corners=5pt, thick,
        minimum width=19cm, align=left, inner sep=6pt,
        font=\sffamily\small, text width=17.5cm},
    infopanel/.style={rectangle, rounded corners=5pt, draw=gray!50, thick,
        minimum width=19cm, align=left, fill=white,
        inner sep=6pt, font=\sffamily\small, text width=17.5cm},
    toolpanel/.style={rectangle, rounded corners=5pt, draw=gray!60, thick,
        minimum width=19cm, align=left, fill=gray!4,
        inner sep=6pt, font=\sffamily\small, text width=17.5cm},
    flowline/.style={->, >=Stealth, very thick, draw=gray!40},
    consumerlabel/.style={font=\sffamily\small\itshape, text=gray!60!black,
        align=center, text width=19cm}
]
    % === STEP BAR ===
    \node[stepbar, draw=gray!60!black, fill=googleGray] (step)
        {\iconSearch{gray!60!black}\;\;
         \textbf{\large GoogleSearcher --- Translational \&
          Clinical Context}
         {\normalsize\itshape\; (optional)}\\[3pt]
         {\small Planner: \textit{``Clinical trials, recent
          developments, therapeutic landscape --- grounds hypotheses
          in translational relevance.''}}\\[2pt]
         {\scriptsize\sffamily\color{gray!60!black}
          Often the last step, but the planner may invoke it
          earlier, omit it entirely, or call it multiple times
          depending on the query.}};

    % === READS ===
    \node[rwpanel, draw=green!40!black, fill=agentGreen,
          below=0.4cm of step] (reads)
        {\textbf{$\blacktriangledown$ Reads from AgentState}\\[5pt]
         $\bullet$\; \texttt{query} {\scriptsize (str)} ---
          research question (distilled to 2--4 keywords by the agent)\\[2pt]
         $\bullet$\; \texttt{previous\_results} {\scriptsize (str)} ---
          full accumulated output including hypotheses from
          ScientistsAgent};

    \draw[flowline] (step.south) -- (reads.north);

    % === PROMPT HIGHLIGHTS ===
    \node[infopanel, below=0.4cm of reads] (prompt)
        {\textbf{System Prompt Highlights}
         {\scriptsize (from \texttt{GoogleSearcher.system\_message})}\\[5pt]
         \textbf{1.}\; ``Extract the essential concepts to form a
          concise, keyword-based search phrase (typically 2--4
          words)''\\[2pt]
         \textbf{2.}\; ``Synthesise report ($\sim$2000 words) ---
          combine facts from all relevant summaries into a single,
          well-structured answer''\\[2pt]
         \textbf{3.}\; ``Highlight points of consensus across sources,
          unique or important details, and any notable
          discrepancies''\\[2pt]
         \textbf{4.}\; ``If summaries do not fully answer the query,
          state what information is missing and suggest a revised
          search phrase''\\[2pt]
         \textbf{5.}\; ``Cite sources using the source title or URL''};

    \draw[flowline] (reads.south) -- (prompt.north);

    % === TOOL ===
    \node[toolpanel, below=0.4cm of prompt] (tool)
        {\textbf{Tool:}\;
         \texttt{google\_search\_tool\_wrapper(query)}\\[5pt]
         {\scriptsize \textbf{Pipeline:}\;
          Google Custom Search API $\rightarrow$
          w3m full-page text extraction $\rightarrow$
          LLM summarisation per page $\rightarrow$
          cross-source synthesis}\\[3pt]
         {\scriptsize \textbf{Returns:}\;
          search results with per-page summaries and source URLs}\\[3pt]
         {\scriptsize \textbf{Resources:}\;
          Google Custom Search Engine\; $\cdot$\;
          w3m text browser\; $\cdot$\;
          LLM relevance filter}\\[3pt]
         {\scriptsize \textit{See companion illustration for the
          Google/web-search source pipeline.}}};

    \draw[flowline] (prompt.south) -- (tool.north);

    % === WRITES ===
    \node[rwpanel, draw=blue!30!black, fill=stateBlue,
          below=0.4cm of tool] (writes)
        {\textbf{$\blacktriangle$ Writes to AgentState}\\[5pt]
         $\bullet$\; \texttt{agent\_outputs[]} {\scriptsize (str)} ---
          synthesised web-research report ($\sim$2000 words),
          including:\\[2pt]
         \hphantom{$\bullet$\;}
          {\scriptsize --- clinical trial status for proposed
           targets (ClinicalTrials.gov)}\\[1pt]
         \hphantom{$\bullet$\;}
          {\scriptsize --- FDA drug labels and regulatory status}\\[1pt]
         \hphantom{$\bullet$\;}
          {\scriptsize --- biotech press releases and pre-print
           developments}\\[1pt]
         \hphantom{$\bullet$\;}
          {\scriptsize --- translational viability assessment with
           source citations}};

    \draw[flowline] (tool.south) -- (writes.north);

    % === DOWNSTREAM CONSUMERS ===
    \node[consumerlabel, below=0.3cm of writes] (consumers)
        {$\downarrow$\; Consumed by:\;
         \textbf{ReportingNode}
          (clinical/translational context, web evidence)};

\end{tikzpicture}%
}% end resizebox
\caption{\textbf{GoogleSearcher} --- workflow interface view.
This agent is typically invoked late in the pipeline---often after
ScientistsAgent has produced mechanistic hypotheses---but the planner
may invoke it at any point, omit it for purely molecular queries,
or re-invoke it after a replan cycle.  Its system prompt defines a
disciplined three-step workflow: (1)~distil the query
into 2--4 core keywords, (2)~execute a targeted Google Custom Search,
and (3)~synthesise the retrieved page summaries into a cohesive
$\sim$2000-word report.  Unlike the PubMed agent, its value lies in
surfacing information \emph{not} indexed in biomedical databases---clinical
trial registries, FDA labels, biotech press releases, and
pre-prints---thereby grounding the proposed targets in translational
reality (see companion illustration for the web-search pipeline).}
\label{fig:closeup-google}
\end{figure}

% ============================================================
\begin{figure}[htbp]
\centering
\resizebox{\textwidth}{!}{%
\begin{tikzpicture}[node distance=0.35cm, font=\sffamily,
    stepbar/.style={rectangle, rounded corners=6pt, thick,
        minimum width=19cm, align=center, font=\sffamily,
        drop shadow={opacity=0.10}, inner sep=6pt, text width=18cm},
    rwpanel/.style={rectangle, rounded corners=5pt, thick,
        minimum width=19cm, align=left, inner sep=6pt,
        font=\sffamily\small, text width=17.5cm},
    infopanel/.style={rectangle, rounded corners=5pt, draw=gray!50, thick,
        minimum width=19cm, align=left, fill=white,
        inner sep=6pt, font=\sffamily\small, text width=17.5cm},
    toolpanel/.style={rectangle, rounded corners=5pt, draw=gray!60, thick,
        minimum width=19cm, align=left, fill=gray!4,
        inner sep=6pt, font=\sffamily\small, text width=17.5cm},
    flowline/.style={->, >=Stealth, very thick, draw=gray!40},
    consumerlabel/.style={font=\sffamily\small\itshape, text=gray!60!black,
        align=center, text width=19cm}
]
    % === STEP BAR ===
    \node[stepbar, draw=gray!60!black, fill=googleGray] (step)
        {\iconSearch{gray!60!black}\;\;
         \textbf{\large GoogleSearcher --- Translational \&
          Clinical Context}
         {\normalsize\itshape\; (optional)}\\[3pt]
         {\small Planner: \textit{``Clinical trials, recent
          developments, therapeutic landscape --- grounds hypotheses
          in translational relevance.''}}\\[2pt]
         {\scriptsize\sffamily\color{gray!60!black}
          Often the last step, but the planner may invoke it
          earlier, omit it entirely, or call it multiple times
          depending on the query.}};

    % === READS ===
    \node[rwpanel, draw=green!40!black, fill=agentGreen,
          below=0.4cm of step] (reads)
        {\textbf{$\blacktriangledown$ Reads from AgentState}\\[5pt]
         $\bullet$\; \texttt{query} {\scriptsize (str)} ---
          research question (distilled to 2--4 keywords by the agent)\\[2pt]
         $\bullet$\; \texttt{previous\_results} {\scriptsize (str)} ---
          full accumulated output including hypotheses from
          ScientistsAgent};

    \draw[flowline] (step.south) -- (reads.north);

    % === PROMPT HIGHLIGHTS ===
    \node[infopanel, below=0.4cm of reads] (prompt)
        {\textbf{System Prompt Highlights}
         {\scriptsize (from \texttt{GoogleSearcher.system\_message})}\\[5pt]
         \textbf{1.}\; ``Extract the essential concepts to form a
          concise, keyword-based search phrase (typically 2--4
          words)''\\[2pt]
         \textbf{2.}\; ``Synthesise report ($\sim$2000 words) ---
          combine facts from all relevant summaries into a single,
          well-structured answer''\\[2pt]
         \textbf{3.}\; ``Highlight points of consensus across sources,
          unique or important details, and any notable
          discrepancies''\\[2pt]
         \textbf{4.}\; ``If summaries do not fully answer the query,
          state what information is missing and suggest a revised
          search phrase''\\[2pt]
         \textbf{5.}\; ``Cite sources using the source title or URL''};

    \draw[flowline] (reads.south) -- (prompt.north);

    % === TOOL ===
    \node[toolpanel, below=0.4cm of prompt] (tool)
        {\textbf{Tool:}\;
         \texttt{google\_search\_tool\_wrapper(query)}\\[5pt]
         {\scriptsize \textbf{Pipeline:}\;
          Google Custom Search API $\rightarrow$
          w3m full-page text extraction $\rightarrow$
          LLM summarisation per page $\rightarrow$
          cross-source synthesis}\\[3pt]
         {\scriptsize \textbf{Returns:}\;
          search results with per-page summaries and source URLs}\\[3pt]
         {\scriptsize \textbf{Resources:}\;
          Google Custom Search Engine\; $\cdot$\;
          w3m text browser\; $\cdot$\;
          LLM relevance filter}\\[3pt]
         {\scriptsize \textit{See companion illustration for the
          Google/web-search source pipeline.}}};

    \draw[flowline] (prompt.south) -- (tool.north);

    % === WRITES ===
    \node[rwpanel, draw=blue!30!black, fill=stateBlue,
          below=0.4cm of tool] (writes)
        {\textbf{$\blacktriangle$ Writes to AgentState}\\[5pt]
         $\bullet$\; \texttt{agent\_outputs[]} {\scriptsize (str)} ---
          synthesised web-research report ($\sim$2000 words),
          including:\\[2pt]
         \hphantom{$\bullet$\;}
          {\scriptsize --- clinical trial status for proposed
           targets (ClinicalTrials.gov)}\\[1pt]
         \hphantom{$\bullet$\;}
          {\scriptsize --- FDA drug labels and regulatory status}\\[1pt]
         \hphantom{$\bullet$\;}
          {\scriptsize --- biotech press releases and pre-print
           developments}\\[1pt]
         \hphantom{$\bullet$\;}
          {\scriptsize --- translational viability assessment with
           source citations}};

    \draw[flowline] (tool.south) -- (writes.north);

    % === DOWNSTREAM CONSUMERS ===
    \node[consumerlabel, below=0.3cm of writes] (consumers)
        {$\downarrow$\; Consumed by:\;
         \textbf{ReportingNode}
          (clinical/translational context, web evidence)};

\end{tikzpicture}%
}% end resizebox
\caption{\textbf{GoogleSearcher} --- workflow interface view.
This agent is typically invoked late in the pipeline---often after
ScientistsAgent has produced mechanistic hypotheses---but the planner
may invoke it at any point, omit it for purely molecular queries,
or re-invoke it after a replan cycle.  Its system prompt defines a
disciplined three-step workflow: (1)~distil the query
into 2--4 core keywords, (2)~execute a targeted Google Custom Search,
and (3)~synthesise the retrieved page summaries into a cohesive
$\sim$2000-word report.  Unlike the PubMed agent, its value lies in
surfacing information \emph{not} indexed in biomedical databases---clinical
trial registries, FDA labels, biotech press releases, and
pre-prints---thereby grounding the proposed targets in translational
reality (see companion illustration for the web-search pipeline).}
\label{fig:closeup-google}
\end{figure}

% ============================================================
\begin{figure}[htbp]
\centering
\resizebox{\textwidth}{!}{%
\begin{tikzpicture}[node distance=0.35cm, font=\sffamily,
    stepbar/.style={rectangle, rounded corners=6pt, thick,
        minimum width=19cm, align=center, font=\sffamily,
        drop shadow={opacity=0.10}, inner sep=6pt, text width=18cm},
    rwpanel/.style={rectangle, rounded corners=5pt, thick,
        minimum width=19cm, align=left, inner sep=6pt,
        font=\sffamily\small, text width=17.5cm},
    infopanel/.style={rectangle, rounded corners=5pt, draw=gray!50, thick,
        minimum width=19cm, align=left, fill=white,
        inner sep=6pt, font=\sffamily\small, text width=17.5cm},
    toolpanel/.style={rectangle, rounded corners=5pt, draw=gray!60, thick,
        minimum width=19cm, align=left, fill=gray!4,
        inner sep=6pt, font=\sffamily\small, text width=17.5cm},
    flowline/.style={->, >=Stealth, very thick, draw=gray!40},
    consumerlabel/.style={font=\sffamily\small\itshape, text=gray!60!black,
        align=center, text width=19cm}
]
    % === STEP BAR ===
    \node[stepbar, draw=gray!60!black, fill=googleGray] (step)
        {\iconSearch{gray!60!black}\;\;
         \textbf{\large GoogleSearcher --- Translational \&
          Clinical Context}
         {\normalsize\itshape\; (optional)}\\[3pt]
         {\small Planner: \textit{``Clinical trials, recent
          developments, therapeutic landscape --- grounds hypotheses
          in translational relevance.''}}\\[2pt]
         {\scriptsize\sffamily\color{gray!60!black}
          Often the last step, but the planner may invoke it
          earlier, omit it entirely, or call it multiple times
          depending on the query.}};

    % === READS ===
    \node[rwpanel, draw=green!40!black, fill=agentGreen,
          below=0.4cm of step] (reads)
        {\textbf{$\blacktriangledown$ Reads from AgentState}\\[5pt]
         $\bullet$\; \texttt{query} {\scriptsize (str)} ---
          research question (distilled to 2--4 keywords by the agent)\\[2pt]
         $\bullet$\; \texttt{previous\_results} {\scriptsize (str)} ---
          full accumulated output including hypotheses from
          ScientistsAgent};

    \draw[flowline] (step.south) -- (reads.north);

    % === PROMPT HIGHLIGHTS ===
    \node[infopanel, below=0.4cm of reads] (prompt)
        {\textbf{System Prompt Highlights}
         {\scriptsize (from \texttt{GoogleSearcher.system\_message})}\\[5pt]
         \textbf{1.}\; ``Extract the essential concepts to form a
          concise, keyword-based search phrase (typically 2--4
          words)''\\[2pt]
         \textbf{2.}\; ``Synthesise report ($\sim$2000 words) ---
          combine facts from all relevant summaries into a single,
          well-structured answer''\\[2pt]
         \textbf{3.}\; ``Highlight points of consensus across sources,
          unique or important details, and any notable
          discrepancies''\\[2pt]
         \textbf{4.}\; ``If summaries do not fully answer the query,
          state what information is missing and suggest a revised
          search phrase''\\[2pt]
         \textbf{5.}\; ``Cite sources using the source title or URL''};

    \draw[flowline] (reads.south) -- (prompt.north);

    % === TOOL ===
    \node[toolpanel, below=0.4cm of prompt] (tool)
        {\textbf{Tool:}\;
         \texttt{google\_search\_tool\_wrapper(query)}\\[5pt]
         {\scriptsize \textbf{Pipeline:}\;
          Google Custom Search API $\rightarrow$
          w3m full-page text extraction $\rightarrow$
          LLM summarisation per page $\rightarrow$
          cross-source synthesis}\\[3pt]
         {\scriptsize \textbf{Returns:}\;
          search results with per-page summaries and source URLs}\\[3pt]
         {\scriptsize \textbf{Resources:}\;
          Google Custom Search Engine\; $\cdot$\;
          w3m text browser\; $\cdot$\;
          LLM relevance filter}\\[3pt]
         {\scriptsize \textit{See companion illustration for the
          Google/web-search source pipeline.}}};

    \draw[flowline] (prompt.south) -- (tool.north);

    % === WRITES ===
    \node[rwpanel, draw=blue!30!black, fill=stateBlue,
          below=0.4cm of tool] (writes)
        {\textbf{$\blacktriangle$ Writes to AgentState}\\[5pt]
         $\bullet$\; \texttt{agent\_outputs[]} {\scriptsize (str)} ---
          synthesised web-research report ($\sim$2000 words),
          including:\\[2pt]
         \hphantom{$\bullet$\;}
          {\scriptsize --- clinical trial status for proposed
           targets (ClinicalTrials.gov)}\\[1pt]
         \hphantom{$\bullet$\;}
          {\scriptsize --- FDA drug labels and regulatory status}\\[1pt]
         \hphantom{$\bullet$\;}
          {\scriptsize --- biotech press releases and pre-print
           developments}\\[1pt]
         \hphantom{$\bullet$\;}
          {\scriptsize --- translational viability assessment with
           source citations}};

    \draw[flowline] (tool.south) -- (writes.north);

    % === DOWNSTREAM CONSUMERS ===
    \node[consumerlabel, below=0.3cm of writes] (consumers)
        {$\downarrow$\; Consumed by:\;
         \textbf{ReportingNode}
          (clinical/translational context, web evidence)};

\end{tikzpicture}%
}% end resizebox
\caption{\textbf{GoogleSearcher} --- workflow interface view.
This agent is typically invoked late in the pipeline---often after
ScientistsAgent has produced mechanistic hypotheses---but the planner
may invoke it at any point, omit it for purely molecular queries,
or re-invoke it after a replan cycle.  Its system prompt defines a
disciplined three-step workflow: (1)~distil the query
into 2--4 core keywords, (2)~execute a targeted Google Custom Search,
and (3)~synthesise the retrieved page summaries into a cohesive
$\sim$2000-word report.  Unlike the PubMed agent, its value lies in
surfacing information \emph{not} indexed in biomedical databases---clinical
trial registries, FDA labels, biotech press releases, and
pre-prints---thereby grounding the proposed targets in translational
reality (see companion illustration for the web-search pipeline).}
\label{fig:closeup-google}
\end{figure}

% ============================================================
\begin{figure}[htbp]
\centering
\resizebox{\textwidth}{!}{%
\begin{tikzpicture}[node distance=0.35cm, font=\sffamily,
    stepbar/.style={rectangle, rounded corners=6pt, thick,
        minimum width=19cm, align=center, font=\sffamily,
        drop shadow={opacity=0.10}, inner sep=6pt, text width=18cm},
    rwpanel/.style={rectangle, rounded corners=5pt, thick,
        minimum width=19cm, align=left, inner sep=6pt,
        font=\sffamily\small, text width=17.5cm},
    infopanel/.style={rectangle, rounded corners=5pt, draw=gray!50, thick,
        minimum width=19cm, align=left, fill=white,
        inner sep=6pt, font=\sffamily\small, text width=17.5cm},
    toolpanel/.style={rectangle, rounded corners=5pt, draw=gray!60, thick,
        minimum width=19cm, align=left, fill=gray!4,
        inner sep=6pt, font=\sffamily\small, text width=17.5cm},
    flowline/.style={->, >=Stealth, very thick, draw=gray!40},
    consumerlabel/.style={font=\sffamily\small\itshape, text=gray!60!black,
        align=center, text width=19cm}
]
    % === STEP BAR ===
    \node[stepbar, draw=gray!60!black, fill=googleGray] (step)
        {\iconSearch{gray!60!black}\;\;
         \textbf{\large GoogleSearcher --- Translational \&
          Clinical Context}
         {\normalsize\itshape\; (optional)}\\[3pt]
         {\small Planner: \textit{``Clinical trials, recent
          developments, therapeutic landscape --- grounds hypotheses
          in translational relevance.''}}\\[2pt]
         {\scriptsize\sffamily\color{gray!60!black}
          Often the last step, but the planner may invoke it
          earlier, omit it entirely, or call it multiple times
          depending on the query.}};

    % === READS ===
    \node[rwpanel, draw=green!40!black, fill=agentGreen,
          below=0.4cm of step] (reads)
        {\textbf{$\blacktriangledown$ Reads from AgentState}\\[5pt]
         $\bullet$\; \texttt{query} {\scriptsize (str)} ---
          research question (distilled to 2--4 keywords by the agent)\\[2pt]
         $\bullet$\; \texttt{previous\_results} {\scriptsize (str)} ---
          full accumulated output including hypotheses from
          ScientistsAgent};

    \draw[flowline] (step.south) -- (reads.north);

    % === PROMPT HIGHLIGHTS ===
    \node[infopanel, below=0.4cm of reads] (prompt)
        {\textbf{System Prompt Highlights}
         {\scriptsize (from \texttt{GoogleSearcher.system\_message})}\\[5pt]
         \textbf{1.}\; ``Extract the essential concepts to form a
          concise, keyword-based search phrase (typically 2--4
          words)''\\[2pt]
         \textbf{2.}\; ``Synthesise report ($\sim$2000 words) ---
          combine facts from all relevant summaries into a single,
          well-structured answer''\\[2pt]
         \textbf{3.}\; ``Highlight points of consensus across sources,
          unique or important details, and any notable
          discrepancies''\\[2pt]
         \textbf{4.}\; ``If summaries do not fully answer the query,
          state what information is missing and suggest a revised
          search phrase''\\[2pt]
         \textbf{5.}\; ``Cite sources using the source title or URL''};

    \draw[flowline] (reads.south) -- (prompt.north);

    % === TOOL ===
    \node[toolpanel, below=0.4cm of prompt] (tool)
        {\textbf{Tool:}\;
         \texttt{google\_search\_tool\_wrapper(query)}\\[5pt]
         {\scriptsize \textbf{Pipeline:}\;
          Google Custom Search API $\rightarrow$
          w3m full-page text extraction $\rightarrow$
          LLM summarisation per page $\rightarrow$
          cross-source synthesis}\\[3pt]
         {\scriptsize \textbf{Returns:}\;
          search results with per-page summaries and source URLs}\\[3pt]
         {\scriptsize \textbf{Resources:}\;
          Google Custom Search Engine\; $\cdot$\;
          w3m text browser\; $\cdot$\;
          LLM relevance filter}\\[3pt]
         {\scriptsize \textit{See companion illustration for the
          Google/web-search source pipeline.}}};

    \draw[flowline] (prompt.south) -- (tool.north);

    % === WRITES ===
    \node[rwpanel, draw=blue!30!black, fill=stateBlue,
          below=0.4cm of tool] (writes)
        {\textbf{$\blacktriangle$ Writes to AgentState}\\[5pt]
         $\bullet$\; \texttt{agent\_outputs[]} {\scriptsize (str)} ---
          synthesised web-research report ($\sim$2000 words),
          including:\\[2pt]
         \hphantom{$\bullet$\;}
          {\scriptsize --- clinical trial status for proposed
           targets (ClinicalTrials.gov)}\\[1pt]
         \hphantom{$\bullet$\;}
          {\scriptsize --- FDA drug labels and regulatory status}\\[1pt]
         \hphantom{$\bullet$\;}
          {\scriptsize --- biotech press releases and pre-print
           developments}\\[1pt]
         \hphantom{$\bullet$\;}
          {\scriptsize --- translational viability assessment with
           source citations}};

    \draw[flowline] (tool.south) -- (writes.north);

    % === DOWNSTREAM CONSUMERS ===
    \node[consumerlabel, below=0.3cm of writes] (consumers)
        {$\downarrow$\; Consumed by:\;
         \textbf{ReportingNode}
          (clinical/translational context, web evidence)};

\end{tikzpicture}%
}% end resizebox
\caption{\textbf{GoogleSearcher} --- workflow interface view.
This agent is typically invoked late in the pipeline---often after
ScientistsAgent has produced mechanistic hypotheses---but the planner
may invoke it at any point, omit it for purely molecular queries,
or re-invoke it after a replan cycle.  Its system prompt defines a
disciplined three-step workflow: (1)~distil the query
into 2--4 core keywords, (2)~execute a targeted Google Custom Search,
and (3)~synthesise the retrieved page summaries into a cohesive
$\sim$2000-word report.  Unlike the PubMed agent, its value lies in
surfacing information \emph{not} indexed in biomedical databases---clinical
trial registries, FDA labels, biotech press releases, and
pre-prints---thereby grounding the proposed targets in translational
reality (see companion illustration for the web-search pipeline).}
\label{fig:closeup-google}
\end{figure}
